\section{Multi-operator responses in the medium}\label{app:D.30}

The rate of a velocity-dependent coupling needs more rows of \cref{eq:8.1-rows} than the density. We write it with those rows, reduce them to a minimal set, relate the set to the published crystal responses and screen each row; \cref{sec:10.4} uses each result. The target is an insulator at $T = 0$ in its ground state $|0\rangle$, with a Hamiltonian whose potentials are local, without spin--orbit coupling and without magnetic order, unless stated otherwise. All energies deposited are positive, $\omega > 0$.

\subsection{The contraction for one velocity-dependent operator}

Take $\mathcal{O}_8 = \vect{S}_\chi\cdot\vect{v}^\perp$ with coefficient $c_8 \equiv c_8^{(e)}(q)$ for dark matter of spin $\frac12$. Its reduction \cref{eq:8.1-O7O8} gives the potentials $\tilde V_{n_e} = c_8\,\vect{S}_\chi\cdot\bar{\vect{v}}_\chi$ and $\tilde V_{j^\alpha_e} = -c_8S^\alpha_\chi$ on the density and the three components of the current. With $\Tr_\chi(S^\alpha_\chi S^\beta_\chi)/2 = \delta_{\alpha\beta}/4$, the tensor \cref{eq:8.1-D} has the elements $\mathfrak{D}_{nn} = c_8^2|\bar{\vect{v}}_\chi|^2/4$, $\mathfrak{D}_{nj^\alpha} = \mathfrak{D}_{j^\alpha n} = -c_8^2\bar v^\alpha_\chi/4$ and $\mathfrak{D}_{j^\alpha j^\beta} = c_8^2\delta_{\alpha\beta}/4$. The contraction of \cref{eq:8.2-rate} is then
\begin{equation}
  \sum_{a,b}\mathfrak{D}_{ab}S_{ab} = \frac{c_8^2}{4}\Bigl[|\bar{\vect{v}}_\chi|^2S_{nn} - 2\Real\bigl(\bar{\vect{v}}_\chi\cdot\vect{S}_{n\vect{j}}\bigr) + \sum_\alpha S_{j^\alpha j^\alpha}\Bigr] ,
  \label{eq:D.30-O8}
\end{equation}
where $\vect{S}_{n\vect{j}}$ is the vector of elements $S_{nj^\alpha}$ and we used $S_{j^\alpha n} = S_{nj^\alpha}^*$, which holds because the response matrix is Hermitian (\cref{eq:4.2-positivity}).

Now split every vector into its component along $\hat{\vect{q}}$ and its components across. With two unit vectors $\hat{\vect{e}}_1$, $\hat{\vect{e}}_2$ perpendicular to $\hat{\vect{q}}$ and to each other, write $j_{\mathrm{L}} = \hat{\vect{q}}\cdot\vect{j}_e(\vect{q})$ and $j_{\mathrm{T}i} = \hat{\vect{e}}_i\cdot\vect{j}_e(\vect{q})$, and likewise $\bar v_{\mathrm{L}} = \hat{\vect{q}}\cdot\bar{\vect{v}}_\chi$ and $\bar{\vect{v}}_{\mathrm{T}} = \bar{\vect{v}}_\chi - \bar v_{\mathrm{L}}\hat{\vect{q}}$. For a local Hamiltonian the continuity equation \cref{eq:4.3-continuity} gives, between eigenstates, $\langle f|j_{\mathrm{L}}|0\rangle = (\omega/q)\langle f|n_e|0\rangle$ with $\omega = E_f - E_0$. Every element of the response matrix that carries $j_{\mathrm{L}}$ is a multiple of one without it:
\begin{equation}
  S_{nj_{\mathrm{L}}} = \frac{\omega}{q}S_{nn} , \qquad S_{j_{\mathrm{L}}j_{\mathrm{L}}} = \frac{\omega^2}{q^2}S_{nn} , \qquad S_{j_{\mathrm{L}}j_{\mathrm{T}i}} = \frac{\omega}{q}S_{nj_{\mathrm{T}i}} .
  \label{eq:D.30-continuity}
\end{equation}
On the dark-matter side, energy conservation \cref{eq:1.3-omegaq} gives $\vect{q}\cdot\bar{\vect{v}}_\chi = \omega + \Delta_\chi$, so $\bar v_{\mathrm{L}} = (\omega + \Delta_\chi)/q$ (\cref{eq:8.1-transverse}). Inserting both into \cref{eq:D.30-O8}, the longitudinal terms collect into $(\bar v_{\mathrm{L}} - \omega/q)^2S_{nn} = (\Delta_\chi/q)^2S_{nn}$, and
\begin{equation}
  \sum_{a,b}\mathfrak{D}_{ab}S_{ab} = \frac{c_8^2}{4}\Bigl[\Bigl(|\bar{\vect{v}}_{\mathrm{T}}|^2 + \frac{\Delta_\chi^2}{q^2}\Bigr)S_{nn} - 2\Real\bigl(\bar{\vect{v}}_{\mathrm{T}}\cdot\vect{S}_{n\vect{j}_{\mathrm{T}}}\bigr) + \Tr S_{\vect{j}_{\mathrm{T}}\vect{j}_{\mathrm{T}}}\Bigr] ,
  \label{eq:D.30-O8T}
\end{equation}
with $\Tr S_{\vect{j}_{\mathrm{T}}\vect{j}_{\mathrm{T}}} = S_{j_{\mathrm{T}1}j_{\mathrm{T}1}} + S_{j_{\mathrm{T}2}j_{\mathrm{T}2}}$. For elastic dark matter only the transverse parts survive, on both sides. The same steps apply to $\mathcal{O}_5$, whose target operator $\vect{q}\times\vect{j}_e$ is already transverse, and to the current part $-2\bar{\vect{v}}_\chi\cdot\vect{j}_e$ of $\mathcal{O}_2$.

\subsection{The kinematic kernel}

The rate per unit target mass needs \cref{eq:D.30-O8T} under the integral of \cref{eq:8.2-spectrum}. Suppose the responses have been averaged over the directions of $\vect{q}$, as in \cref{app:A.8}. The averaged cross block $\bar{\vect{S}}_{n\vect{j}_{\mathrm{T}}}$ is then a vector perpendicular to $\hat{\vect{q}}$ that cannot depend on the orientation of $\hat{\vect{e}}_1$ and $\hat{\vect{e}}_2$, since rotating them about $\hat{\vect{q}}$ is part of the average; it vanishes, and the cross term drops out. For the two diagonal terms the velocity integral at fixed $q$ and $\omega$ is done as in \cref{app:A.8}. With $\vect{v}$ at polar angle $\theta$ to $\hat{\vect{q}}$, the delta function fixes $v\cos\theta = \vmin$, so that $|\bar{\vect{v}}_{\mathrm{T}}|^2 = v^2\sin^2\theta = v^2 - \vmin^2$ (the shift $\vect{q}/2\mchi$ of $\bar{\vect{v}}_\chi$ is longitudinal). The angular integral is $\Theta(v - \vmin)/(2qv)$ as before, and per unit $\dd^3q/(2\pi)^3$, for an isotropic $f_\chi$ or after the average over $\hat{\vect{q}}$ of \cref{app:A.8},
\begin{equation}
  \begin{gathered}
  \int\dd^3v\,f_\chi(\vect{v})\,\delta(\omega - \omega_{\vect{q}})\,\bigl\{1,\;|\bar{\vect{v}}_{\mathrm{T}}|^2\bigr\} = \frac{1}{2q}\bigl\{\eta(\vmin),\;\eta_2(\vmin)\bigr\} ,
  \\
  \eta_2(\vmin) = \int\dd^3v\,f_\chi\,\frac{v^2 - \vmin^2}{v}\,\Theta(v - \vmin) ,
  \end{gathered}
  \label{eq:D.30-kernel}
\end{equation}
where the second form of $\eta_2$ equals the definition $2\int_{\vmin}^\infty\dd v'\,v'\eta(v')$ of \cref{app:A.8}: exchanging the order of the integrals, $\int_{\vmin}^{v}2v'\,\dd v' = v^2 - \vmin^2$. As for $\eta$, only the average of the halo over directions enters once the response is averaged.

\subsection{The minimal set of rows}

Three symmetries reduce the rows of \cref{eq:8.1-rows} for the electrons.

\emph{Continuity.} \Cref{eq:D.30-continuity} removes $j_{\mathrm{L}}$. If the Hamiltonian also conserves the spin of each electron, the spin current obeys its own continuity equation, $\partial_ts^\beta_e(\vect{q},t) = \ii q^\alpha\mathcal{J}^{\alpha\beta}_e(\vect{q},t)$, by the same commutator with the kinetic energy that leads to \cref{eq:4.3-continuity}, and $\hat q^\alpha\mathcal{J}^{\alpha\beta}_e$ acts as $(\omega/q)s^\beta_e$. What remains of the currents are the transverse components $j_{\mathrm{T}i}$ and $\mathcal{J}^{\mathrm{T}i,\beta}_e \equiv \hat e^\alpha_i\mathcal{J}^{\alpha\beta}_e$.

\emph{Spin rotations.} Without spin--orbit coupling and without magnetic order, the Hamiltonian commutes with the total spin and the ground state of an insulator with filled bands is a singlet. The density, the current and the kinetic-energy density are scalars under spin rotations; $\vect{s}_e$ and $\mathcal{J}^{\alpha\cdot}_e$ are vectors in their spin index. A scalar acting on the singlet produces singlets, and a vector produces triplets, which are orthogonal to them. Every element of the response matrix between a scalar row and a vector row vanishes, and by the Wigner--Eckart theorem the vector--vector elements are $\delta_{\beta\beta'}$ times one function.

For independent particles the two blocks are related. The determinant of filled bands has both spins of every valence orbital occupied. A singlet row $O$ has the one-electron vertex $o_{cv}\,\delta_{\sigma\sigma'}$ between orbitals with spins $\sigma$ and $\sigma'$; its triplet partner $O\otimes\sigma^z/2$ has $o_{cv}\,(\sigma^z)_{\sigma'\sigma}/2$. Summing the products of two vertices over the four spin combinations gives $\sum_{\sigma\sigma'}\delta_{\sigma\sigma'} = 2$ for the singlet and $\sum_{\sigma\sigma'}|(\sigma^z)_{\sigma'\sigma}|^2/4 = \frac12$ for the triplet, and the spatial sums are the same. The $x$ and $y$ components give $\Tr(\sigma^x\sigma^x)/4 = \frac12$ too. Hence
\begin{equation}
  \begin{gathered}
  S^0_{s^\beta s^{\beta'}} = \frac{\delta_{\beta\beta'}}{4}\,S^0_{nn} , \qquad S^0_{s^\beta\mathcal{J}^{\mathrm{T}i,\beta'}} = \frac{\delta_{\beta\beta'}}{4}\,S^0_{nj_{\mathrm{T}i}} , \\
  S^0_{\mathcal{J}^{\mathrm{T}i,\beta}\mathcal{J}^{\mathrm{T}k,\beta'}} = \frac{\delta_{\beta\beta'}}{4}\,S^0_{j_{\mathrm{T}i}j_{\mathrm{T}k}} ,
  \end{gathered}
  \label{eq:D.30-quarter}
\end{equation}
with the superscript 0 for independent particles. Beyond independent particles the blocks differ, because the Coulomb interaction acts on the total density $n_\uparrow + n_\downarrow$ and not on the spin density $n_\uparrow - n_\downarrow$. The Hartree kernel and its local fields enter the singlet block only. In density-functional language the kernel of the triplet block is the spin part of $f_{\mathrm{xc}}$, $(f_{\uparrow\uparrow} - f_{\uparrow\downarrow})/2$, without $v_{\mathrm{C}}$.

\emph{Inversion with time reversal.} In silicon and germanium the antiunitary product $\mathsf{PT}$ of inversion and time reversal is a symmetry. Time reversal maps $n_e(\vect{q})$ to $n_e(-\vect{q})$ and $\vect{j}_e(\vect{q})$ to $-\vect{j}_e(-\vect{q})$, since it reverses momenta; inversion maps $\vect{q}$ back to $-\vect{q}$ and reverses $\vect{p}$ again. Both rows are even under $\mathsf{PT}$ at the same $\vect{q}$, and so is $\mathcal{T}_e$. An antiunitary symmetry of the ground state and of $H$ maps $\langle0|O_a^\dagger|f\rangle\langle f|O_b|0\rangle$, summed over a degenerate set of final states, to its complex conjugate, and every element of the singlet block is real. The same holds for the triplet block: its rows change sign under $\mathsf{PT}$, since time reversal also flips the spin, and the products of two of them stay real.

\emph{Symmetry axes.} When $\hat{\vect{q}}$ lies on a two-, three- or fourfold axis of the cubic group, the symmetries that leave $\vect{q}$ fixed include a rotation about $\hat{\vect{q}}$ by $\pi$, $2\pi/3$ or $\pi/2$. Such a rotation leaves $n_e(\vect{q})$ unchanged and rotates $(j_{\mathrm{T}1}, j_{\mathrm{T}2})$ as a vector in the transverse plane. A vector in that plane that is unchanged by the rotation is zero, so $S_{nj_{\mathrm{T}i}} = 0$. On the three- and fourfold axes the current block is also a multiple of the unit matrix.

\subsection{Response functions $\overline W_1$--$\overline W_5$}

Catena and collaborators write the general rate with the density vertex $f = \langle i'\vect{k}'|\ee^{\ii\vect{q}\cdot\vect{x}}|i\vect{k}\rangle$ and the vector vertex $\vect{f} = \langle i'\vect{k}'|\ee^{\ii\vect{q}\cdot\vect{x}}\,\ii\nabla/m_e|i\vect{k}\rangle$~\cite[Eqs.~(10), (11)]{catena2021crystal}. They use five bilinears, $B_1 = |f|^2$, $B_2 = (\vect{q}/m_e)\cdot f\vect{f}^*$, $B_3 = |\vect{f}|^2$, $B_4 = |(\vect{q}/m_e)\cdot\vect{f}|^2$ and $B_5 = \ii(\vect{q}/m_e)\cdot(\vect{f}\times\vect{f}^*)$, each summed over pairs into a response $\overline W_l(q,\omega)$~\cite[Eqs.~(37)--(42), (47)]{catena2021crystal}. Since $\ii\nabla = -\vect{p}$, a plane wave $\vect{k}$ gives $\vect{f}$ the element $-\vect{k}/m_e$, while the current vertex of \cref{eq:8.1-rows} is $(\vect{k} + \vect{k}')/2m_e = \vect{k}/m_e + \vect{q}/2m_e$. Between any two states $\vect{f} = (\vect{q}/2m_e)f - \vect{j}$, with $\vect{j}$ the current vertex. The continuity equation reads $\vect{q}\cdot\vect{j} = \omega f$, and
\begin{equation}
  \vect{q}\cdot\vect{f} = -\Bigl(\omega - \frac{q^2}{2m_e}\Bigr)f , \qquad \hat{\vect{e}}_i\cdot\vect{f} = -j_{\mathrm{T}i} .
  \label{eq:D.30-catena}
\end{equation}
Hence $\overline W_2 = -(\omega - q^2/2m_e)\overline W_1/m_e$ and $\overline W_4 = (\omega - q^2/2m_e)^2\overline W_1/m_e^2$; the transverse combination $\overline W_3 - (m_e^2/q^2)\overline W_4$ is the sum of $|j_{\mathrm{T}i}|^2$; and $\overline W_5 = 0$ and $\Imag\overline W_2 = 0$ when the vertices are real, as $\mathsf{PT}$ makes them. Comparing $B_1$ with \cref{eq:10.1-S0}, and using the relation $\overline W_1 = (8\omega\alpha m_e^2/q^3)|f_{\mathrm{crys}}|^2$~\cite[Eq.~(D2)]{catena2021crystal} with \cref{eq:10.1-fcrys}, the dictionary is $\overline W_1 = (2\Omega\omega/\pi)\bar S^0_{nn}$ and, by the same sum over pairs, $\overline W_3 - (m_e^2/q^2)\overline W_4 = (2\Omega\omega/\pi)\Tr\bar S^0_{\vect{j}_{\mathrm{T}}\vect{j}_{\mathrm{T}}}$. A nonlocal pseudopotential breaks \cref{eq:D.30-catena} through the commutator $[V_{\mathrm{NL}}, \ee^{\ii\vect{q}\cdot\vect{r}}]$, and the current vertex then needs the correction \cref{eq:E.7-vNL}.

\subsection{Screening of each row}

The screened response of any rows follows from \cref{eq:8.4-screened}. Work in the reciprocal-lattice basis of \cref{sec:4.6}, with the Coulomb interaction acting on the densities $n_e(\vect{q} + \vect{G})$ and the rows $a$, $b$ entering only as external vertices. The random-phase approximation with local fields gives
\begin{equation}
  R_{ab}(\vect{q},\omega) = \tilde R_{ab} + \sum_{\vect{G}\vect{G}'}\tilde R_{a,\vect{G}}\,\bigl[v_{\mathrm{C}}(1 - \tilde R\,v_{\mathrm{C}})^{-1}\bigr]_{\vect{G}\vect{G}'}\,\tilde R_{\vect{G}',b} , \qquad
  S_{ab} = \ii\bigl(R_{ab} - R_{ba}^*\bigr) ,
  \label{eq:D.30-dyson}
\end{equation}
where $v_{\mathrm{C}}$ is the diagonal matrix $e^2/|\vect{q} + \vect{G}|^2$, $\tilde R_{a,\vect{G}}$ is the proper response of row $a$ to the density at $\vect{q} + \vect{G}$, and $\tilde R = R^0$ at this level. The second relation is the anti-Hermitian part of $R$ at $T = 0$, the generalization of $S = -2\Imag R$ to off-diagonal elements. The correction to row $a$ is set by its proper correlation with the densities. For $a = n_e(\vect{q})$ itself it is the screening of \cref{sec:10.2}, with the head $\vect{G} = 0$ dominant at small $q$. For $a = j_{\mathrm{T}i}$ with $\hat{\vect{q}}$ on a symmetry axis, $\tilde R_{j_{\mathrm{T}i},\vect{0}} = 0$ by the argument of the previous subsection, applied to the proper response, so only $\vect{G} \neq 0$ contributes: the transverse current is screened through local fields alone. For $a = \mathcal{T}_e$, a scalar under the point group, $\tilde R_{\mathcal{T},\vect{0}} \neq 0$ and the head acts. For the triplet rows every $\tilde R_{a,\vect{G}}$ vanishes by spin rotation, and \cref{eq:D.30-dyson} leaves them unscreened.

The transverse block has a measurable limit. At $\vect{q} \to 0$ the dielectric tensor of \cref{sec:4.4} is $\varepsilon^{\alpha\beta} = \delta^{\alpha\beta} - (e^2/\omega^2)[\nd_e\delta^{\alpha\beta}/m_e + \tilde R^{\alpha\beta}_{jj}]$, so that $\Imag\varepsilon^{\mathrm{T}} = -(e^2/\omega^2)\Imag\tilde R_{j_{\mathrm{T}}j_{\mathrm{T}}}$, and with $S = -2\Imag R$ and $R_{j_{\mathrm{T}}j_{\mathrm{T}}} = \tilde R_{j_{\mathrm{T}}j_{\mathrm{T}}}$ on a symmetry axis without local fields,
\begin{equation}
  \Tr S_{\vect{j}_{\mathrm{T}}\vect{j}_{\mathrm{T}}} = \frac{4\omega^2}{e^2}\,\Imag\varepsilon^{\mathrm{T}}(\vect{q},\omega) , \qquad
  S_{nn} = \frac{2q^2}{e^2}\,\Imag\Bigl[-\frac{1}{\varepsilon(\vect{q},\omega)}\Bigr] , \qquad
  \frac{\Tr S_{\vect{j}_{\mathrm{T}}\vect{j}_{\mathrm{T}}}}{S_{nn}} \to \frac{2\omega^2}{q^2}\,|\varepsilon(\omega)|^2 ,
  \label{eq:D.30-optical}
\end{equation}
the last form at $q \to 0$ in a cubic crystal, where $\varepsilon^{\mathrm{T}} \to \varepsilon$ (\cref{eq:4.4-ET}). The second relation is \cref{eq:4.4-SL}. For independent particles the same steps give $2(\omega/q)^2$ without the factor $|\varepsilon|^2$.

